\eabstract{
In recent years in Japan, the working-age population has been declining in many regions due to a falling birthrate, an aging population, and outmigration, making it increasingly difficult to address regional challenges. In response, the Ministry of Education, Culture, Sports, Science and Technology (MEXT) has, in part, placed expectations on regional universities to contribute to solving these issues. In fact, many universities are strengthening research rooted in their local communities in order to demonstrate their own significance.
On the other hand, from a bibliometric perspective, there has been little development of indicators that capture such region-oriented research activities. This is likely due to several challenges. First, it is difficult to accurately extract place names (toponyms) from texts and other materials. Second, it is also challenging to uniquely identify and disambiguate those place names. Ultimately, there will be a need for mechanisms to uniquely link place names and geographic information to the relevant papers and datasets.
While this study takes such solutions into consideration, it first focuses on the accuracy of place name extraction and conducts an experimental study using a compact test dataset.
}
